% ============================================================
% A03 — Produk: Daftar, Tambah, Edit
% ============================================================
\flowmeta{A03}{Produk: Daftar, Tambah, Edit}{Admin}

\begin{flowsection}{Tujuan}
Kelola katalog: daftar produk, tambah, ubah, dan kelola varian serta foto.
\end{flowsection}

\begin{flowsection}{Prasyarat}
Login sebagai admin dengan peran yang memiliki akses modul ini.
\end{flowsection}

\begin{flowsection}{Langkah}
\begin{enumerate}
  \item Buka Produk \&  Varian SKU (/admin/products).
  \item Tambah produk: tombol Tambah (ikon plus) → isi nama, kategori, SKU, harga, berat, foto.
  \item Aktifkan varian bila perlu (atribut bertingkat), lalu simpan.
  \item Ubah produk dari menu titik tiga (MoreVertical) pada baris.
\end{enumerate}
\end{flowsection}

\begin{flowsection}{Catatan Stok}
Form produk \textbf{tidak} menetapkan stok. Produk/varian dibuat dengan
\textbf{stok 0}; stok hanya bertambah melalui \textbf{Purchase Order $\rightarrow$
Penerimaan Barang (GRN)} dan berpindah melalui transfer antar-gudang. Varian
disimpan sebagai data otonom pada tabel \texttt{product\_variants} (id asli
dipakai oleh PO/GRN).
\end{flowsection}

\shot{gambar/admin/a03-1.png}{Halaman Produk: Daftar, Tambah, Edit pada Panel Admin}{fig:a03-1}

\begin{flowsection}{Hasil yang Diharapkan}
Produk tersimpan dan tampil di katalog storefront.
\end{flowsection}

\begin{flowsection}{Penanganan Error dan Kasus Khusus}
\begin{itemize}
  \item Field wajib (nama, kategori, SKU, harga, foto) harus diisi.
\end{itemize}
\end{flowsection}

\begin{flowsection}{Kaitan Modul}
\begin{itemize}
  \item \texttt{ProductListPage.jsx}
  \item \texttt{ProductCreateForm.jsx}
  \item \texttt{ProductController}
  \item \texttt{/api/products}
\end{itemize}
\end{flowsection}

\docver{v1.0}{Sep 2026}
